\subsection{The exact weight-seven image}
\label{subsec:weight-seven-exact}

The uniform odd theorem supplies the integral image; the leading-space
calculation identifies the corresponding full degree-seven kernel.

\begin{theorem}\label{thm:weight-seven-exact-image}
One has
\[
 m_7(D_5)=m_7(J_5)=128\mathbf Z_2.
\]
Consequently, for $D_7=D_5\cap\ker m_7$ and $J_7=D_7\cap H$,
\[
 D_5=D_7J_5,\qquad G=D_7H,\qquad
 B\simeq D_7/\overline{[G,D_7]J_7}.
\]
The subgroup $D_7$ is proper, closed and normal in $G$, and
\[
 \eta_7:D_5\longrightarrow\mathbf F_2,
 \qquad \eta_7(g)=m_7(g)/128\pmod2
\]
is a continuous $G$-invariant character with nonzero arithmetic
restriction.
\end{theorem}

\begin{proof}
Here $D_5=C_2$. Theorem~\ref{thm:uniform-odd-arithmetic-image}
with $r=3$ gives the two exact images and the factorization
$D_5=D_7J_5$, hence $G=D_7H$. The scalar conjugation law makes
$D_7$ closed and normal and makes $m_7/128$ invariant modulo two.
The explicit arithmetic value in
Proposition~\ref{prop:weight-seven-arithmetic-value} proves properness
and nonvanishing on $J_5$. The canonical relative quotient is the
centralization argument of Corollary~\ref{cor:cubic-radical-normalization}.
Finally, Lemma~\ref{lem:weight-seven-integral-line} shows that killing
$m_7$ kills the entire degree-seven leading term.
\end{proof}

\begin{corollary}\label{cor:weight-seven-nonextension}
The character $\eta_7$ has no $G$-invariant continuous extension to
$D_3$, $N_{\mathrm{rad}}$ or $N$.
Moreover, writing $\psi_q=\rho_q^I/2\pmod2$ for the divided
odd-prime Kummer characters on $N$, one still has
\[
 \ker\bigl[H^1(N,\mathbf F_2)^G\longrightarrow
                         H^1(D_7,\mathbf F_2)\bigr]
 =\left(\bigoplus_{q\text{ odd prime}}\mathbf F_2\psi_q\right)
       \oplus\mathbf F_2\eta_3.
\]
On $D_3$, the corresponding kernel is exactly $\mathbf F_2\eta_5$.
\end{corollary}

\begin{proof}
Use the arithmetic elements constructed above, for which
\[
 z_3\in J_3,\qquad
 z_5=h z_3h^{-1}z_3^{-t^5}\in J_5,
 \qquad \eta_7(z_5)=1.
\]
For every continuous character $\chi$ to $\mathbf F_2$ and every
profinite exponent $c$, continuity of powers gives
$\chi(z^c)=(c\bmod2)\chi(z)$. If $\chi$ is $G$-invariant on
$D_3$, then
\[
 \chi(z_5)=\chi(z_3)-(t_2^5\bmod2)\chi(z_3)=0,
\]
because $t_2$ is odd. Thus its restriction to $D_5$ cannot be
$\eta_7$. Restricting an invariant character from either larger
group to $D_3$ proves the other nonextension assertions.

Now suppose that a $G$-invariant character of $N$ vanishes on
$D_7$. Its restriction to $D_5$ factors through
$D_5/D_7=128\mathbf Z_2$, so is zero or $\eta_7$.
Nonextension excludes the latter, and hence it vanishes on $D_5$.
Its restriction to $D_3$ is similarly zero or $\eta_5$.
The weight-five arithmetic word
$z_3=h z_1h^{-1}z_1^{-t^3}$, with $z_1\in N_{\mathrm{rad}}$,
satisfies $\eta_5(z_3)=1$, whereas the same odd-power calculation
makes every invariant character of $N$ zero on $z_3$.
Consequently the character vanishes on $D_3$.
The established cubic joint quotient
\[
 N/D_3\simeq
 \left(\prod_{q\text{ odd prime}}2\widehat{\mathbf Z}\right)
       \times8\mathbf Z_2
\]
now gives exactly the displayed direct sum. Conversely every one
of those characters already vanishes on $D_3$ and hence on $D_7$.
For a character defined on $D_3$, the first nonextension argument
again forces vanishing on $D_5$; the quotient
$D_3/D_5=32\mathbf Z_2$ then gives precisely $\mathbf F_2\eta_5$.
\end{proof}

If $J_5=G_{L_5}$ and $L_7=\overline{\mathbf Q}^{\,J_7}$, the exact
image also gives $\operatorname{Gal}(L_7/L_5)\simeq\mathbf Z_2$.
This specifies the extension by the weight-seven monodromy
coordinate; no expression by rational radicals is asserted.
